\documentclass[11pt]{article}

\usepackage[margin=1in]{geometry}
\usepackage{array}
\usepackage{booktabs}
\usepackage{enumitem}
\usepackage{hyperref}
\usepackage{microtype}
\usepackage{tabularx}
\usepackage[T1]{fontenc}
\usepackage{lmodern}

\hypersetup{
    colorlinks=true,
    linkcolor=blue,
    urlcolor=blue
}

\newcolumntype{Y}{>{\raggedright\arraybackslash}X}

\title{\textit{cs-mail} Deployment and Migration Profile\\\large Request Payments and the C-SQD Member Program}
\author{C-SQD}
\date{Draft 0.6 --- September 2026}

\setlength{\emergencystretch}{1.5em}

\begin{document}
\maketitle

\begin{abstract}
An initial C-SQD deployment can operate the communication service and its
canonical ledger while external providers execute request-specific charges,
refunds, and member payments. This profile follows those obligations from a
sender's relationship request to a quarterly member distribution. It defines the
C-SQD financial program, then explains the operational, privacy, and migration
choices needed to deliver it. Company ownership of the pool does not make its
member-restricted portion available for ordinary corporate spending.
\end{abstract}

\section{What This Profile Controls}

The \textit{Protocol Specification} defines permission, request outcomes, and
the distinction between committed obligations and completed external payments.
Section~\ref{sec:program}, the \emph{C-SQD Financial Program}, adds normative
requirements for deployments advertising this profile. In that section,
``must'' and ``must not'' are requirements. The remaining sections are deployment
guidance. Values explicitly marked proposed are planning defaults, not approved
live-payment terms.

This revision replaces periodic account funding, direct recipient collateral
payouts, and withdrawals of spendable ledger balances. Its financial narrative
is a conditional service charge for one relationship request, followed by an
outcome-dependent refund or forfeiture and, separately, member distributions.
The legal characterization and payment-provider arrangements still require
assessment; these labels do not establish a money-transmission exemption.

\section{A Request from Charge to Outcome}

Alice has no standing permission to write to Bob. She reviews one request quote,
authorizes $C+S$, and supplies her initial encrypted communication. C-SQD records
one request and asks the external provider to capture that charge under a stable
operation key. Once capture evidence and durable content are available, admission
opens the request and starts its decision clock.

Bob's follow-up setting governs additional messages during that clock. If he
permits them, Alice may add messages under the same request. They cost no more,
do not move the deadline, and do not generate another request prompt. If he does
not permit them, C-SQD refuses the follow-up without selling another charge
against the pending request.

Acceptance records a full refund obligation and standing permission. Rejection
records processing revenue $C$ and pending forfeiture $S$. Silence through expiry
records $C$ as revenue and a refund obligation for $S$. Bob receives no personal
entitlement from rejecting Alice. Each obligation remains visible until the
external execution is reconciled, even if Alice or Bob closes an account.

Capture and admission cannot form one database transaction with a payment
provider. Recovery must therefore complete the story. A cancelled or timed-out
preparing request is not admitted when a delayed capture confirmation arrives;
the captured amount is refunded in full. An uncertain capture is queried under
its original key rather than retried as another charge. A failed refund remains
an obligation and enters retry or review.

\section{C-SQD Financial Program}\label{sec:program}

\subsection{Ownership and the limits on its use}

C-SQD is the service operator and owner of the company pool in this profile.
It is responsible for request refunds and finalized member payments, while an
external provider executes the underlying transactions. The live policy must
identify the contracting entity, merchant/payment responsibilities, and applicable
provider arrangements before charges begin.

The pool has no participant ownership shares and no earmarked claim linking a
particular rejected request to its recipient. Before quarterly allocation, a
member has no individually finalized claim to a fraction of its unallocated
funds. At allocation, the ledger records an individual payment obligation for
each qualifying member. This accounting distinction must be reflected in the
published program terms rather than inferred from the word ``pool.''

Eligible forfeitures undergo one disclosed corporate allocation. The remainder
is restricted to member distributions. C-SQD must not spend that remainder on
operations, payout fees, or corporate losses, or assess it again in later
quarters. Processing revenue and the corporate allocation fund those company
costs. This rule qualifies the earlier proposal to devote all forfeitures to
members; it does not authorize further undisclosed deductions.

\subsection{When a forfeiture becomes eligible}

Rejection is the beginning of the forfeiture process, not permission to distribute
cash. Each pending lot must retain a scoped request cause, original payment
reference, amount, unit, rate/policy version, maturity evidence, and assessment
status. Before eligibility, the policy's capture/settlement, refund, dispute,
and reserve conditions must be satisfied. Refund obligations and unresolved
amounts must remain separately accounted for and unavailable to the pool.

Quarter end alone does not establish payment finality. A lot that fails maturity
at cutoff stays pending until a later eligible quarter; it does not receive an
estimated member allocation. A hold or dispute must identify the affected lot
and its review condition. Reserve methodology, release conditions, and later-loss
funding must be approved before live distribution, not invented by a worker.

C-SQD must record later external reversals through authorized compensating
entries and fund the resulting obligations without silently taking restricted
member funds or cancelling finalized payables. A final request outcome remains
part of the journal even when a financial correction is necessary.

\subsection{The corporate allocation happens once}

The proposed corporate rate is $\alpha=0.03$. The adopted rate and rounding
policy must be disclosed through immutable request terms. A future rate change
applies prospectively; the rate at a later maturity date does not replace the
one attached to an existing request.

At quarter $q$, group newly eligible, previously unassessed lots by settlement
unit and quoted rate version $v$. Let $F_{qv}$ be a group's amount in integer
minor units. The corporate amount and new member contribution are
\[
G_q=\sum_v\left\lfloor\alpha_v F_{qv}\right\rfloor,
\qquad M_q=\sum_v F_{qv}-G_q.
\]
The floor is applied once per rate cohort and unit at the finalized quarter
boundary, with any fractional corporate unit left for members. Each contributing
lot must be marked assessed in the same atomic allocation record. Replays cannot
assess it again. Calculations must use checked integer or exact rational
arithmetic, never floating-point currency.

For example, $10{,}000$ units of newly eligible forfeitures at 3\% produce $300$
of corporate revenue and $9{,}700$ for members. This assessment excludes $C$,
refund obligations, unresolved request amounts, and any refundable persistence
reserve. Previously assessed member funds and finalized unpaid allocations are
also excluded. The sender is not charged an additional 3\%.

\subsection{Who qualifies for an equal share}

The program benefits qualifying active members generally. It must not weight a
share by requests received, collateral contributed, rejections, messages sent,
or correspondence with a particular sender. One qualifying member receives one
share, irrespective of aliases, mailboxes, devices, or subscription seats.

The distribution identity is distinct from the protocol's private principal and
public communication identities. The program must prevent duplicate enrollment
while keeping identity evidence out of delivery, relationship, and telemetry
records. Eligibility requires program opt-in, the published membership/activity
conditions, and required identity checks. Confirmed abuse or duplicate membership
must be handled under a disclosed review and correction procedure.

The proposed initial eligibility policy is membership for at least 30 days by
cutoff and intentional authenticated use on at least two distinct days of the
quarter. Reading can count. Background synchronization, automated polling, and
incoming lane traffic alone do not establish activity. Store minimal activity
evidence, not message content or a history of correspondents. Whether legal
entities or managed groups qualify independently remains unresolved; mail
enrollment does not itself settle that question.

\subsection{From quarter cutoff to fixed payables}

Allocations occur quarterly. The program must publish its calendar, timezone,
eligibility policy version, and cutoff rule before the period. A calendar quarter
is represented as a half-open interval ending at the next quarter's start; an
event at that boundary belongs to the new quarter. UTC is the proposed reference
timezone. A review interval after cutoff can resolve already recorded evidence
without admitting later activity into the closed quarter.

At finalization, fix the qualifying membership snapshot and the eligible funding
manifest. Let $U_{q-1}$ be unallocated, previously assessed member carryforward,
and let $N_q$ be the qualifying member count. The amount available for equal
allocation is $D_q=U_{q-1}+M_q$. For $N_q>0$,
\[
a_q=\left\lfloor D_q/N_q\right\rfloor,
\qquad U_q=D_q-N_q a_q.
\]
Each member receives a payable of $a_q$ minor units. The remainder $U_q$ remains
restricted carryforward. If $N_q=0$, no individual allocation is created and
$U_q=D_q$. An empty pool creates no payment. Settlement units must be handled
separately; no implicit foreign exchange is permitted.

The snapshot, manifest, assessment marks, payables, remainder, and journal entries
must finalize as one visible outcome, using a prepared batch if necessary.
An operator must not silently rerun the quarter with a different membership set.
Corrections require a separate authorized record with a stated cause and funding
source; they do not overwrite the finalized batch.

\subsection{Paying a member}

Finalized payables are obligations, not available pool cash. A member with a
failed payment, incomplete payout onboarding, or an amount below a disclosed
minimum threshold retains that payable. It must not return to the next quarter's
equal split. Several fixed payables may be paid together through one provider
operation while preserving their individual accounting causes.

Payment requires a stable operation key and verified destination under the
external provider's rules. The obligation is discharged only by reconciled
payment confirmation. Retries, account closure, and quarter rollover do not
erase it. The proposed operational target is payment initiation within 30 days
after quarter end; this is distinct from a promise of provider completion.

Member statements must separate finalized allocation, awaiting payment, confirmed
payment, and any disclosed threshold carryforward. There is no spendable cs-mail
balance and no transfer to another member. C-SQD must not deduct payout fees from
the stated member allocation under this profile.

\subsection{Published policy and program closure}

The launch policy must resolve the proposed rate, eligibility thresholds, legal-
entity eligibility, maturity/reserve criteria, payment thresholds, and review
timelines. It must also define account closure, unclaimed obligations, residual
restricted funds, and orderly program termination before collecting real value.
Stopping the program cannot by itself convert member-restricted funds into
unrestricted corporate revenue or extinguish unpaid obligations.

The program should publish aggregate reconciliation: new eligible forfeitures,
corporate allocation, member contribution, carryforward, qualifying member count,
finalized payables, and payment completion. Individual correspondence or the
amount a particular recipient rejected must not be disclosed through those
reports.

\section{Operating the Service Without a Wallet}

C-SQD can initially run one provider and ledger, but identity, relationship,
ciphertext, accounting, membership verification, and telemetry need distinct
credentials, keys, and access roles. The payment adapter receives scoped payment
references and required amounts. It does not need message content, a recipient's
contact graph, or events for every follow-up message.

The operational ledger should reconcile captures, unresolved requests, refunds,
pending forfeitures, corporate revenue, restricted funds, and member payables
against external statements. An operator investigates mismatches through narrowly
authorized records and compensating commands, not arbitrary balance edits.
Unknown outcomes stay visible until resolved. A processor outage must not become
a duplicate charge or a claim that a refund has completed.

Communication enrollment may support individuals, organizations, managed groups,
and limited accounts. Collect control and recovery evidence appropriate to each
role; do not demand legal identity merely to populate an ordinary mailbox.
Distribution verification is a separate purpose with its own access, retention,
and eligibility rules. Provider migration should preserve recipient-scoped abuse
continuity without exporting a universal principal identifier.

\section{The Recipient and Sender Experience}

The sender sees the price of one relationship request, its initial communication,
the possible refunds, and the decision window. After submission, the interface
shows preparing, awaiting decision, or the terminal outcome, with external refund
status separately. It does not present available or withdrawable funds. A
preparing request offers authenticated cancellation; closing the application
does not imply that the charge was cancelled.

The recipient sees one grouped request and the choices accept, reject, and block.
The explanation starts with future permission, then the financial consequence.
Rejecting does not earn the recipient that request's collateral. Dismissing the
prompt is not a decision. A late acceptance grants future permission without
reopening the closed request's refund outcome.

Follow-up controls belong beside request preferences and are enforced at
admission. Recipients may disable follow-ups or set bounded counts and rates;
changes affect subsequent admissions, not immutable financial terms. Accepted
communication and valid express lanes retain their own authorization paths.

The proposed request policy uses a 30-day decision window, a 30-day cooldown
after expiry, and a 90-day cooldown after rejection. Cooldowns run from the
canonical decision deadline for expiry and the canonical timely rejection/block
event for rejection allocation, subject to any later principal-history bound.
These values remain proposals. A finalized policy must define recipient-imposed
longer limits, explicit invitations that can shorten ordinary cooldowns, and
behavior after revocation. Blocking always requires an authenticated unblock;
an invitation cannot silently remove it. The prior nonzero persistence reserve
is not enabled in the base draft and requires a separate decision and extension.

\section{SMTP and Express-Lane Migration}

SMTP is a visibly identified compatibility route because its gateway may handle
plaintext. A native message's end-to-end confidentiality must not be claimed for
ordinary gateway-readable mail.

Inbound processing first checks blocking, standing acceptance, and an applicable
recipient-granted lane. It then checks whether an authenticated message may join
an identified open request. Only an eligible new request reaches the payment
invitation flow. The web authorization binds the directed parties, request,
initial content/declarations, amount, immutable terms, and deadlines. It never
asks the sender to top up a general account.

Subsequent legacy mail must prove its association and satisfy current follow-up
policy; an address or subject-line match alone is insufficient authority. A
failed lane or follow-up cannot silently authorize a new charge. Invitations
need rate limits, anti-enumeration controls, clear expiry, and abuse reporting.
Outbound downgrade must be disclosed before the sender commits to that route.

Express-lane messages carry no request charge or pool contribution. Granting or
ending a lane does not settle a pending request. Domain-authentication evidence
supports the lane's sender scope, not the truth of its declared purpose or a
member's active-use eligibility.

\section{Incentives, Review, and Later Federation}

Equal member distributions remove a direct rejection bounty, but they do not
remove economic incentives. C-SQD earns $C$ on unsuccessful requests and a share
of eligible forfeitures. It should publish pricing controls and preserve signed
decision evidence so neither the service nor a recipient can manufacture expiry
by hiding acceptance. Review should address duplicate members, fabricated
activity, abusive request patterns, collusion, and manipulated cutoff evidence.

A single company-owned pool is the initial model. Independent-provider federation
requires a new agreement assigning the merchant, refund debtor, pool owner,
membership authority, duplicate-member controls, clearing, default losses, and
payment responsibility. Netting may discharge valid obligations more efficiently;
it cannot create new ownership rules or turn a rejected sender's collateral into
a direct recipient payout. The centralized draft therefore does not promise that
the former recipient-collateral clearing example transfers unchanged to federation.

\section{A Staged Launch}

Start with a deterministic model of one request and its outcomes. Add durable
ordering, double-entry obligations, and recovery through a simulated external
provider. A desktop pilot can use nonredeemable simulated amounts while testing
real encryption, request/follow-up permission, deadlines, and failure behavior.

Before live value, the service needs approved payment and program policies,
external reconciliation, capture-after-cancellation recovery, refund retries,
maturity controls, reproducible quarterly snapshots, and failed-payout handling.
The first distribution should be demonstrable end to end with duplicate provider
events and a failed member payment before real members depend on it. Privacy
tests must cover verification, activity evidence, payment records, and backups as
well as messages.

SMTP and independent-provider federation are later operational gates. Each gate
needs measurable exit evidence appropriate to its risks; a working communications
pilot is not evidence that live financial operations or cross-provider pool
governance are complete.

\end{document}
